% T4 fig:barrier (opus1) — Eyring activation landscape, equilibrium vs driven, as the causal split.
% Reaction coordinate x in [-1.5,1.5]; wells at x=-1 (filled site), x=+1 (empty site + ion), TS at x=0.
% G_eq(x)=dGa cos^2(pi x/2), dGa=1;  G_drv(x)=G_eq - A(x+1)/2, A=0.6, chi=1/2  (coords: T4_calc.py)
% => forward barrier dGa-chi*A = 0.70, reverse dGa+(1-chi)A = 1.30, product well down by A = 0.60.
\centering
\begin{tikzpicture}[x=1.75cm,y=2.05cm,font=\scriptsize]
% ================= (a) equilibrium =================
\begin{scope}
\draw[-{Latex[length=1.4mm]}] (-1.62,-0.72) -- (1.70,-0.72);
\node[below,font=\scriptsize] at (0.04,-0.80) {reaction coordinate};
\draw[-{Latex[length=1.4mm]}] (-1.62,-0.72) -- (-1.62,1.18) node[above,font=\scriptsize] {free energy};
\draw[thick] plot[smooth] coordinates {(-1.500,0.5000) (-1.350,0.2730) (-1.200,0.0955) (-1.050,0.0062) (-0.900,0.0245) (-0.750,0.1464) (-0.600,0.3455) (-0.450,0.5782) (-0.300,0.7939) (-0.150,0.9455) (0.000,1.0000) (0.150,0.9455) (0.300,0.7939) (0.450,0.5782) (0.600,0.3455) (0.750,0.1464) (0.900,0.0245) (1.050,0.0062) (1.200,0.0955) (1.350,0.2730) (1.500,0.5000)};
\node[below,font=\scriptsize] at (-1.0,-0.02) {filled site};
\node[below,font=\scriptsize] at (1.0,-0.02) {empty $+$ ion};
\node[above,font=\scriptsize] at (0,1.0) {transition state};
\draw[densely dotted,gray] (-1.0,0) -- (-0.5,0);
\draw[densely dotted,gray] (0,1.0) -- (-0.5,1.0);
\draw[{Latex[length=1.2mm]}-{Latex[length=1.2mm]}] (-0.5,0.03) -- (-0.5,0.97) node[pos=0.5,left,font=\scriptsize] {$\Delta G_a$};
\node[font=\scriptsize] at (0,-1.02) {(a) equilibrium: $\mathcal A=0$};
\end{scope}
% ================= (b) driven =================
\begin{scope}[xshift=5.7cm]
\draw[-{Latex[length=1.4mm]}] (-1.62,-0.72) -- (1.70,-0.72);
\node[below,font=\scriptsize] at (0.04,-0.80) {reaction coordinate};
\draw[densely dotted,thin,gray] plot[smooth] coordinates {(-1.500,0.5000) (-1.200,0.0955) (-0.900,0.0245) (-0.600,0.3455) (-0.300,0.7939) (0.000,1.0000) (0.300,0.7939) (0.600,0.3455) (0.900,0.0245) (1.200,0.0955) (1.500,0.5000)};
\draw[thick] plot[smooth] coordinates {(-1.500,0.6500) (-1.350,0.3780) (-1.200,0.1555) (-1.050,0.0212) (-0.900,-0.0055) (-0.750,0.0714) (-0.600,0.2255) (-0.450,0.4132) (-0.300,0.5839) (-0.150,0.6905) (0.000,0.7000) (0.150,0.6005) (0.300,0.4039) (0.450,0.1432) (0.600,-0.1345) (0.750,-0.3786) (0.900,-0.5455) (1.050,-0.6088) (1.200,-0.5645) (1.350,-0.4320) (1.500,-0.2500)};
\draw[-{Latex[length=1.3mm]},thick] (0,1.0) -- (0,0.72);
\node[above,font=\scriptsize] at (0.05,1.0) {peak $\downarrow\chi\mathcal A$};
% forward barrier dGa - chi A
\draw[densely dotted,gray] (-1.0,0) -- (-0.44,0);
\draw[densely dotted,gray] (0,0.70) -- (-0.44,0.70);
\draw[{Latex[length=1.2mm]}-{Latex[length=1.2mm]}] (-0.44,0.03) -- (-0.44,0.67) node[pos=0.5,left,font=\scriptsize,align=right] {$\Delta G_a$\\$-\chi\mathcal A$};
% reverse barrier dGa + (1-chi) A
\draw[densely dotted,gray] (1.0,-0.60) -- (0.48,-0.60);
\draw[densely dotted,gray] (0,0.70) -- (0.48,0.70);
\draw[{Latex[length=1.2mm]}-{Latex[length=1.2mm]}] (0.48,-0.57) -- (0.48,0.67) node[pos=0.60,right,font=\scriptsize,align=left] {$\Delta G_a$\\$+(1{-}\chi)\mathcal A$};
% affinity A
\draw[densely dotted,gray] (-1.0,0) -- (1.34,0);
\draw[{Latex[length=1.2mm]}-{Latex[length=1.2mm]}] (1.34,0) -- (1.34,-0.60) node[pos=0.5,right,font=\scriptsize] {$\mathcal A$};
\node[below,font=\scriptsize] at (-1.02,-0.02) {filled};
\node[below,font=\scriptsize] at (0.86,-0.62) {empty $+$ ion};
\node[font=\scriptsize] at (0,-1.02) {(b) driven $\mathcal A>0$ ($\chi=\tfrac12$)};
\end{scope}
\end{tikzpicture}
\caption{활성화 장벽 --- 평형과 구동의 인과 분해. (a) 평형($\mathcal A{=}0$) --- 두 골(찬 자리 / 빈 자리$+$이온)이
같은 높이, 장벽 $\Delta G_a$ 위로 넘는 속도가 Eyring 식 $k\simeq k_0e^{-\Delta G_a/RT}$. (b) 구동력 $\mathcal A_j{=}sF(V_n-U_j)$
가 도착 골을 $\mathcal A$ 만큼 내려($\chi{=}\tfrac12$ 예) 정상점을 $\chi\mathcal A$ 낮추면, 정방향 장벽 $\Delta G_a-\chi\mathcal A$,
역방향 $\Delta G_a+(1-\chi)\mathcal A$ 로 비대칭 이동한다(식~\eqref{eq:bv}; 두 속도의 비가 식~\eqref{eq:db} detailed
balance --- $\chi$ 소거). 점선은 (a)의 평형 개형. 온도는 분모 $RT$ 로 진입. 좌표는 $\Delta G_a\cos^2(\pi x/2)$ 에 선형
tilt 를 더한 수치 평가($\Delta G_a{=}1,\ \mathcal A{=}0.6$ 임의 단위). 이 그림이 식~\eqref{eq:xieq} logistic 과 \S\ref{sec:lag}
의 $L_q$ 의 공통 출발점이다.}
\label{fig:barrier}
